\begin{lemma}
Let $\Omega\subseteq\mathbb{C}$ be open, let $u\colon\mathbb{C}\to\mathbb{C}$
and let $D_0,D_1\colon\mathbb{C}\to\mathbb{C}$ satisfy
\begin{enumerate}
\item[(H1)] $u\in L^{1}_{\mathrm{loc}}(\Omega)$ and
  $D_j\in L^{1}_{\mathrm{loc}}(\Omega)$ for $j=0,1$, in the sense of
  \noderef{MemLpLocOn};
\item[(H2)] $D_j$ is the distributional $j$-th partial derivative of $u$ on
  $\Omega$ for $j=0,1$, in the sense of \noderef{IsWeakPartialOn};
\item[(H3)] $u$ is differentiable at almost every point of $\Omega$, in the
  real Fr\'echet sense.
\end{enumerate}
Then for $j=0,1$,
\[
  D_j(z)=\partial_j u(z)\qquad\text{for almost every }z\in\Omega,
\]
where $\partial_0$ and $\partial_1$ are the classical partial derivatives
along $1$ and $i$, that is $\partial_ju=D^{[j]}u$ in the sense of
\noderef{PartialWord}.

The tablet carries two unrelated notions of derivative: $W^{k,s}_{\mathrm{loc}}$
and quasiregularity are written with the integration-by-parts identity of
\noderef{IsWeakPartialOn}, while the Wirtinger operators, the polar derivatives
of Section 4 and the iterated partial derivatives $D^\nu$ of Sections 5 and 6
are all built from the real Fr\'echet derivative.  This lemma is the bridge.
The hypotheses are exactly what the data of a $W^{1,1}_{\mathrm{loc}}$
function provides, with $D_j$ the member of the coherent family indexed by the
one-letter word $[j]$, together with the almost-everywhere differentiability
that item (C2) of the paper's note supplies for the inverse of a quasiconformal
homeomorphism.  Neither (H2) nor (H3) is dispensable.  Without (H3) the
classical derivative need not exist at any point, so there is nothing to
compare $D_j$ with.  And (H3) alone does not make the classical derivative a
distributional one: the Cantor function on $[0,1]$, extended to the plane so as
not to depend on the second coordinate, is continuous and differentiable almost
everywhere with classical derivative $0$ almost everywhere, yet $0$ is not its
distributional derivative.  So the proof below must consume (H2), and it does,
at Step 2.
\end{lemma}
\begin{proof}
Write $d\sigma$ for plane Lebesgue measure, $e_0=1$ and $e_1=i$, so that
$\partial_j v(z)=\mathrm{d}v(z)[e_j]$ and, by the recursion of
\noderef{PartialWord}, $D^{[j]}v=\partial_jv$.  Put $w_j:=\partial_ju$, defined
at every point where $u$ is differentiable.  Fix $j\in\{0,1\}$.

\emph{Step 0: reduction to a compact set.}  Every open subset of the plane is a
countable union of compact sets, say $\Omega=\bigcup_{m\ge1}K^{(m)}$ with each
$K^{(m)}\subseteq\Omega$ compact.  A countable union of null sets is null, so
it suffices to prove that $D_j=w_j$ almost everywhere on each $K^{(m)}$.  Fix
therefore a compact $K\subseteq\Omega$.  Since $K$ is compact, $\Omega$ is open
and $K\subseteq\Omega$, there is $\delta>0$ with
\[
  K^{3\delta}:=\{z:\operatorname{dist}(z,K)\le3\delta\}\subseteq\Omega ;
\]
indeed either $\Omega=\mathbb{C}$, where any $\delta$ serves, or
$\operatorname{dist}(K,\mathbb{C}-\Omega)>0$ because $K$ is compact and
$\mathbb{C}-\Omega$ is closed and disjoint from $K$, and then any $\delta$
below a third of that distance serves.  The set $K^{3\delta}$ is closed and
bounded, hence compact, and $K\subseteq K^{\delta}\subseteq K^{2\delta}
\subseteq K^{3\delta}$.

\emph{Step 1: the mollifier.}  Fix once and for all a function
$\varrho\colon\mathbb{C}\to\mathbb{R}$ of class $C^\infty$ with
$\varrho\ge0$, $\varrho(y)=0$ for $|y|\ge1$ and $\int\varrho\,d\sigma=1$, and
for $0<\varepsilon\le\delta$ put
$\varrho_\varepsilon(y)=\varepsilon^{-2}\varrho(y/\varepsilon)$, which is
$C^\infty$, vanishes for $|y|\ge\varepsilon$ and has
$\int\varrho_\varepsilon\,d\sigma=1$.  Let
\[
  U=\mathbf{1}_{K^{3\delta}}\,u,\qquad V_j=\mathbf{1}_{K^{3\delta}}\,D_j .
\]
By (H1) applied to the compact set $K^{3\delta}\subseteq\Omega$ and
\noderef{MemLpLocOn}, $U$ and $V_j$ belong to $L^{1}(\mathbb{C})$.  Put
$u_\varepsilon=U*\varrho_\varepsilon$ and
$v_{j,\varepsilon}=V_j*\varrho_\varepsilon$.

\emph{Step 2: $u_\varepsilon$ is smooth and
$\partial_ju_\varepsilon=v_{j,\varepsilon}$ near $K$.}  The convolution of an
$L^{1}$ function with a compactly supported $C^\infty$ function is $C^\infty$,
with derivatives obtained by differentiating the smooth factor; so
$u_\varepsilon$ is $C^\infty$ on $\mathbb{C}$ and
\begin{equation}\tag{1}
  \partial_ju_\varepsilon(z)
  =\bigl(U*\partial_j\varrho_\varepsilon\bigr)(z)
  =\int U(w)\,(\partial_j\varrho_\varepsilon)(z-w)\,d\sigma(w).
\end{equation}
Fix $z$ with $\operatorname{dist}(z,K)\le2\delta$ and $0<\varepsilon\le\delta$,
and define $\varphi(w)=\varrho_\varepsilon(z-w)$, read as a
complex-valued function.  Then $\varphi$ is $C^\infty$ and vanishes off the
closed ball $\{w:|w-z|\le\varepsilon\}$, which is compact and contained in
$K^{3\delta}\subseteq\Omega$ because
$\operatorname{dist}(z,K)\le2\delta$ and $\varepsilon\le\delta$.  Hence
$\operatorname{tsupp}\varphi\subseteq\Omega$ and $\varphi$ is a test function
for $\Omega$ in the sense of \noderef{IsTestFn}.  The map $w\mapsto z-w$ is
affine with derivative $-\mathrm{id}$, so by the chain rule
\[
  \mathrm{d}\varphi(w)[e_j]=-(\partial_j\varrho_\varepsilon)(z-w).
\]
Therefore (H2) for this $j$, applied to $\varphi$, reads
\[
  -\int u(w)\,(\partial_j\varrho_\varepsilon)(z-w)\,d\sigma(w)
  =-\int D_j(w)\,\varrho_\varepsilon(z-w)\,d\sigma(w),
\]
that is
\begin{equation}\tag{2}
  \int u(w)\,(\partial_j\varrho_\varepsilon)(z-w)\,d\sigma(w)
  =\int D_j(w)\,\varrho_\varepsilon(z-w)\,d\sigma(w).
\end{equation}
Both integrands vanish off $\{w:|w-z|\le\varepsilon\}\subseteq K^{3\delta}$,
where $u=U$ and $D_j=V_j$; so by (1) the left side of (2) is
$\partial_ju_\varepsilon(z)$ and the right side is
$v_{j,\varepsilon}(z)$.  Hence
\begin{equation}\tag{3}
  \partial_ju_\varepsilon(z)=v_{j,\varepsilon}(z)
  \qquad\bigl(\operatorname{dist}(z,K)\le2\delta,\ 0<\varepsilon\le\delta\bigr).
\end{equation}

\emph{Step 3: the finite-difference identity for $u_\varepsilon$.}  Let
$0<|t|\le\delta$, let $z\in K$ and let $s$ lie between $0$ and $t$.  Then
$\operatorname{dist}(z+se_j,K)\le|s|\le\delta\le2\delta$, so (3) applies at
$z+se_j$.  The function $s\mapsto u_\varepsilon(z+se_j)$ is $C^\infty$ on
$\mathbb{R}$ with derivative $\partial_ju_\varepsilon(z+se_j)$, by the chain
rule, so the fundamental theorem of calculus and (3) give
\begin{equation}\tag{4}
  u_\varepsilon(z+te_j)-u_\varepsilon(z)
  =\int_0^t\partial_ju_\varepsilon(z+se_j)\,ds
  =\int_0^tv_{j,\varepsilon}(z+se_j)\,ds .
\end{equation}

\emph{Step 4: letting $\varepsilon\to0$ in (4).}  For $g\in L^{1}(\mathbb{C})$
one has $\|g*\varrho_\varepsilon-g\|_{L^{1}(\mathbb{C})}\to0$ as
$\varepsilon\to0$; applied to $U$ and $V_j$,
\begin{equation}\tag{5}
  \|u_\varepsilon-U\|_{L^{1}(\mathbb{C})}\to0,\qquad
  \|v_{j,\varepsilon}-V_j\|_{L^{1}(\mathbb{C})}\to0 .
\end{equation}
Fix $t$ with $0<|t|\le\delta$ and regard both sides of (4) as functions of
$z\in K$.  For the left side, translation invariance of $d\sigma$ gives
\[
  \int_K\bigl|u_\varepsilon(z+te_j)-U(z+te_j)\bigr|\,d\sigma(z)
  \le\|u_\varepsilon-U\|_{L^{1}(\mathbb{C})},
\]
and the same bound without the translation, so
$z\mapsto u_\varepsilon(z+te_j)-u_\varepsilon(z)$ converges in $L^{1}(K)$ to
$z\mapsto U(z+te_j)-U(z)$, which equals $u(z+te_j)-u(z)$ for $z\in K$ because
$z$ and $z+te_j$ both lie in $K^{3\delta}$.  For the right side, write $I_t$
for the interval between $0$ and $t$; the integrand is measurable on
$K\times I_t$, and by Tonelli's theorem and translation invariance
\[
  \int_K\Bigl|\int_0^t\bigl(v_{j,\varepsilon}-V_j\bigr)(z+se_j)\,ds\Bigr|
    d\sigma(z)
  \le\int_{I_t}\!\Bigl(\int_K\bigl|v_{j,\varepsilon}-V_j\bigr|(z+se_j)
    \,d\sigma(z)\Bigr)ds
  \le|t|\,\|v_{j,\varepsilon}-V_j\|_{L^{1}(\mathbb{C})},
\]
which tends to $0$ by (5); and $V_j(z+se_j)=D_j(z+se_j)$ for $z\in K$ and
$|s|\le\delta$.  Two $L^{1}(K)$ limits of one and the same family agree almost
everywhere, so
\begin{equation}\tag{6}
  u(z+te_j)-u(z)=\int_0^tD_j(z+se_j)\,ds
  \qquad\text{for a.e. }z\in K,
\end{equation}
for each fixed $t$ with $0<|t|\le\delta$.

\emph{Step 5: the difference quotients converge in $L^{1}(K)$.}  For
$0<|t|\le\delta$ put $\Delta_t(z)=\bigl(u(z+te_j)-u(z)\bigr)/t$.  By (6) and
$\frac1t\int_0^t D_j(z)\,ds=D_j(z)$,
\[
  \Delta_t(z)-D_j(z)=\frac1t\int_0^t\bigl[D_j(z+se_j)-D_j(z)\bigr]ds
  \qquad\text{for a.e. }z\in K,
\]
so, by Tonelli's theorem again, by $D_j=V_j$ on $K^{3\delta}$ and by
$K\subseteq K^{3\delta}$, $K+se_j\subseteq K^{3\delta}$ for $|s|\le\delta$,
\[
  \int_K\bigl|\Delta_t-D_j\bigr|\,d\sigma
  \le\frac1{|t|}\int_{I_t}\Bigl(\int_{\mathbb{C}}
      \bigl|V_j(z+se_j)-V_j(z)\bigr|\,d\sigma(z)\Bigr)ds
  \le\sup_{|s|\le|t|}\bigl\|V_j(\cdot+se_j)-V_j\bigr\|_{L^{1}(\mathbb{C})} ,
\]
the last step because the inner integral is at most the supremum and
$|I_t|=|t|$.
Translation is continuous in $L^{1}(\mathbb{C})$, so the right-hand side tends
to $0$ as $t\to0$.  Hence
\begin{equation}\tag{7}
  \Delta_t\to D_j\quad\text{in }L^{1}(K)\ \text{as }t\to0 .
\end{equation}

\emph{Step 6: the pointwise limit.}  By (H3) there is a null set $N$ such that
$u$ is differentiable at every $z\in\Omega-N$.  For such a $z$ the definition
of the real Fr\'echet derivative gives, for real $t\to0$,
\[
  \Delta_t(z)=\frac{u(z+te_j)-u(z)}{t}\longrightarrow
  \mathrm{d}u(z)[e_j]=w_j(z).
\]

\emph{Step 7: conclusion on $K$.}  Put $t_l=\delta/l$ for $l\ge1$.  By (7),
$\Delta_{t_l}\to D_j$ in $L^{1}(K)$, so some subsequence
$\Delta_{t_{l_i}}$ converges to $D_j$ almost everywhere on $K$.  By Step 6 the
same subsequence converges to $w_j(z)$ at every $z\in K-N$.  A sequence of
complex numbers has at most one limit, so $D_j=w_j$ almost everywhere on $K$.

\emph{Step 8: conclusion on $\Omega$.}  By Step 0 the set
$\{z\in\Omega:D_j(z)\neq w_j(z)\}$ is contained in a countable union of null
sets, one for each $K^{(m)}$, hence is null.  As $j\in\{0,1\}$ was arbitrary,
this is the assertion.
\end{proof}
